\documentclass[letterpaper,10pt]{moderncv}

% moderncv options
\moderncvstyle{casual} % Choose a style: 'casual', 'classic', 'banking', 'oldstyle', 'fancy'
\moderncvcolor{blue} % Choose a color: 'blue', 'orange', 'green', 'red', 'purple', 'grey'

% Personal information
\name{Anvesh Varma}{Dantuluri}
\address{New York, USA}{} % Assuming New York is current location for simplicity
\phone[mobile]{+1 (201) 640-9986}
\email{ad7647@nyu.edu}
\social[linkedin]{linkedin.com/in/anvesh-varma}
\social[github]{github.com/Anveshvarmad}

% Adjust margins for ATS compatibility if needed (defaults are usually good)
% \usepackage[scale=0.75]{geometry} % Default in moderncv

% Custom commands for bullet points
\usepackage{enumitem}
\setlist[itemize]{
  label=\textbullet,
  leftmargin=*,
  itemsep=0em, % Spacing between items
  parsep=0em,  % Spacing between paragraphs within an item
  partopsep=0em, % Spacing above the first item
  topsep=0.2em, % Spacing above the entire list
}

% --- Custom commands for section formatting ---
% Redefine \section to use \section* to suppress numbering and customize spacing
\makeatletter
\renewcommand{\section}[1]{%
  \par\addvspace{1.5ex}%
  \phantomsection{}%
  \addcontentsline{toc}{section}{#1}%
  {\raggedright\Large\itshape\bfseries\upshape\color{color2}#1\par}% % Using color2 for headings
  \nobreak\addvspace{0.5ex}\@afterheading%
}
\makeatother

% Adjust section title color to match moderncvcolor
\colorlet{color2}{black} % Fallback if moderncvcolor isn't defined yet

\begin{document}

\maketitle % Prints the name and contact information

\section{Education}
\cventry{Aug 2024 – May 2026}{Master of Science, Computer Engineering}{\textbf{New York University}}{New York, USA}{}{
  \begin{itemize}
    \item Key Coursework: \textbf{Machine Learning}, \textbf{Deep Learning}, \textbf{Reinforcement Learning}, \textbf{Big Data}, Advanced Python for Data Science, Computer System Architecture, Applied Matrix Theory, Systems Engineering
    \item Migrated legacy Python CGI scripts to a modern \textbf{Django} web application, deployed using Apache + mod\_wsgi, improving maintainability and runtime efficiency
    \item Certifications: \textbf{Data Analysis with Python} (IBM), Python Data Structures (University of Michigan), Google Data Analytics (Google), Data Visualization with Tableau (UC Davis)
  \end{itemize}
}

\section{Technical Skills}
\begin{itemize}
    \item \textbf{Programming \& Data:} Python, SQL, C++, Bash
    \item \textbf{Machine Learning:} Supervised Learning, Unsupervised Learning, Feature Engineering, Model Evaluation, \textbf{AI}
    \item \textbf{Deep Learning:} PyTorch, TensorFlow, Neural Networks, CNNs, Transfer Learning
    \item \textbf{LLMs \& Generative AI:} Large Language Models, Prompt Engineering, RAG, Embeddings
    \item \textbf{MLOps \& Deployment:} Model Versioning, Experiment Tracking, Docker, CI/CD, Inference Optimization
    \item \textbf{Cloud \& AI Systems:} AWS (EC2, S3, Lambda, Bedrock), GPU Computing, Cost Optimization
\end{itemize}

\section{Work Experience}
\cventry{May 2025 – Aug 2025}{Software Development Engineer Intern (\textbf{AI Systems})}{\textbf{INFOLOB}}{Remote}{}{
  \begin{itemize}
    \item Designed a fault-tolerant \textbf{AI orchestration layer} on MCP servers integrating \textbf{AWS Bedrock} for \textbf{inference}
    \item Implemented autonomous fallback and recovery mechanisms increasing task completion by 33\% and availability by 40\%
    \item Optimized end-to-end \textbf{inference pipelines} sustaining sub-minute latency under high-concurrency workloads
    \item Reduced operational cost by 28\% through intelligent model routing and dynamic endpoint selection
  \end{itemize}
}
\cventry{May 2024 – Jul 2024}{Software Engineering Intern (\textbf{Data \& Analytics})}{\textbf{INFOTECH LABS}}{Hyderabad, India}{}{
  \begin{itemize}
    \item Developed scalable \textbf{Power BI dashboards} using \textbf{SQL} and DAX accelerating insight delivery by 20\%
    \item Analyzed high-volume sales datasets identifying trends driving
  \end{itemize}
}

\end{document}